\chapter{Introduction and Foundations}

\section{Introduction to Dynamical Systems}
Linear Systems Theory forms the mathematical bedrock for modern control engineering, signal processing, and systems biology. Dynamic physical systems are characterized by variables that evolve over time. These variables can be categorized into:
\begin{itemize}
    \item \textbf{Inputs} $u(t) \in \R^m$: Exogenous signals applied to the system (e.g., voltages, actuator forces).
    \item \textbf{Outputs} $y(t) \in \R^p$: Measured responses of the system (e.g., positions, temperatures).
    \item \textbf{States} $x(t) \in \R^n$: The minimal set of internal variables that, together with the future input signals, completely determine the future behavior of the system.
\end{itemize}

A system is classified as \textit{linear} if it satisfies the principles of superposition and homogeneity. It is \textit{time-invariant} if its behavior is independent of the absolute initial time. In this textbook, we primarily study Linear Time-Invariant (LTI) systems, represented in state-space form.

\begin{mydefinition}{Continuous-Time LTI State-Space Model}
A continuous-time Linear Time-Invariant (LTI) state-space representation is governed by the system of equations:
\begin{align}
    \dot{x}(t) &= A x(t) + B u(t) \label{eq:lti_state} \\
    y(t) &= C x(t) + D u(t) .\label{eq:lti_output}
\end{align}
where:
\begin{itemize}
    \item $x(t) \in \R^n$ is the state vector.
    \item $u(t) \in \R^m$ is the input vector.
    \item $y(t) \in \R^p$ is the output vector.
    \item $A \in \R^{n \times n}$ is the dynamic (state) matrix, describing the internal couplings.
    \item $B \in \R^{n \times m}$ is the input coupling matrix.
    \item $C \in \R^{p \times n}$ is the output measurement matrix.
    \item $D \in \R^{p \times m}$ is the direct transmission (feedthrough) matrix.
\end{itemize}
\end{mydefinition}

\begin{mydefinition}{Discrete-Time LTI State-Space Model}
In digital implementations, systems are governed by discrete-time dynamics represented by:
\begin{align}
    x_{k+1} &= A_d x_k + B_d u_k \label{eq:lti_discrete_state} \\
    y_k &= C x_k + D u_k .\label{eq:lti_discrete_output}
\end{align}
where $k \in \mathbb{Z}_{\ge 0}$ denotes the discrete sample index, $A_d$ is the discrete dynamic matrix, and $B_d$ is the discrete input matrix.
\end{mydefinition}

\section{Linearization of Multi-Input Multi-Output (MIMO) Systems}
Most physical systems are fundamentally nonlinear. However, local analysis and controller design can be performed by linearizing the nonlinear dynamics about a desired equilibrium point.

Consider an autonomous, continuous-time MIMO nonlinear system described by:
\begin{align}
    \dot{x}(t) &= f(x(t), u(t)) \\
y(t) &= h(x(t), u(t)) .
\end{align}
where $f: \R^n \times \R^m \to \R^n$ and $h: \R^n \times \R^m \to \R^p$ are continuously differentiable vector-valued functions.

Let $(x_e, u_e)$ be a chosen operating (equilibrium) point satisfying the algebraic relation:
\begin{equation}
f(x_e, u_e) = 0 .
\end{equation}
We define the small deviation variables representing local perturbations about this equilibrium point:
\begin{equation}
\delta x(t) = x(t) - x_e, \quad \delta u(t) = u(t) - u_e, \quad \delta y(t) = y(t) - h(x_e, u_e) .
\end{equation}
Applying a multivariate Taylor series expansion to $f(x, u)$ and $h(x, u)$ about $(x_e, u_e)$ yields:
\begin{align}
    f(x, u) &\approx f(x_e, u_e) + \left. \frac{\partial f}{\partial x} \right|_{(x_e, u_e)} (x - x_e) + \left. \frac{\partial f}{\partial u} \right|_{(x_e, u_e)} (u - u_e) + \text{H.O.T.} \\
h(x, u) &\approx h(x_e, u_e) + \left. \frac{\partial h}{\partial x} \right|_{(x_e, u_e)} (x - x_e) + \left. \frac{\partial h}{\partial u} \right|_{(x_e, u_e)} (u - u_e) + \text{H.O.T.} .
\end{align}
Neglecting the Higher-Order Terms (H.O.T.) and substituting the deviation variables gives:
\begin{align}
    \delta \dot{x}(t) &= A \delta x(t) + B \delta u(t) \\
\delta y(t) &= C \delta x(t) + D \delta u(t) .
\end{align}
where $A, B, C, D$ are the Jacobian matrices evaluated at the equilibrium point:
\begin{align}
    A &= \left. \frac{\partial f}{\partial x} \right|_{(x_e, u_e)} = \begin{bmatrix}
        \frac{\partial f_1}{\partial x_1} & \cdots & \frac{\partial f_1}{\partial x_n} \\
        \vdots & \ddots & \vdots \\
        \frac{\partial f_n}{\partial x_1} & \cdots & \frac{\partial f_n}{\partial x_n}
    \end{bmatrix}_{(x_e, u_e)} \\
    B &= \left. \frac{\partial f}{\partial u} \right|_{(x_e, u_e)} = \begin{bmatrix}
        \frac{\partial f_1}{\partial u_1} & \cdots & \frac{\partial f_1}{\partial u_m} \\
        \vdots & \ddots & \vdots \\
        \frac{\partial f_n}{\partial u_1} & \cdots & \frac{\partial f_n}{\partial u_m}
    \end{bmatrix}_{(x_e, u_e)} \\
    C &= \left. \frac{\partial h}{\partial x} \right|_{(x_e, u_e)} = \begin{bmatrix}
        \frac{\partial h_1}{\partial x_1} & \cdots & \frac{\partial h_1}{\partial x_n} \\
        \vdots & \ddots & \vdots \\
        \frac{\partial h_p}{\partial x_1} & \cdots & \frac{\partial h_p}{\partial x_n}
    \end{bmatrix}_{(x_e, u_e)} \\
    D &= \left. \frac{\partial h}{\partial u} \right|_{(x_e, u_e)} = \begin{bmatrix}
        \frac{\partial h_1}{\partial u_1} & \cdots & \frac{\partial h_1}{\partial u_m} \\
        \vdots & \ddots & \vdots \\
        \frac{\partial h_p}{\partial u_1} & \cdots & \frac{\partial h_p}{\partial u_m}
\end{bmatrix}_{(x_e, u_e)} .
\end{align}
This linearization enables engineers to apply the vast body of linear systems tools to locally stabilize and control inherently nonlinear systems.

\section{Linear Algebra Review}
Working with state-space models requires a solid grasp of matrix operations and matrix analysis.

\subsection{Matrix Multiplications}
Consider matrices $A \in \R^{n \times m}$ and $B \in \R^{m \times p}$. The standard product $C = AB \in \R^{n \times p}$ can be viewed through two key structural interpretations:
\begin{enumerate}
    \item \textbf{Inner Product Expansion}: The $(i, j)$-th entry of $C$ is the inner product of the $i$-th row of $A$ and the $j$-th column of $B$:
    \begin{equation}
c_{ij} = \langle l_i^A, c_j^B \rangle = \sum_{k=1}^m a_{ik} b_{kj} .
    \end{equation}
    \item \textbf{Outer Product Expansion}: The product $AB$ can be represented as a sum of rank-one matrices constructed by multiplying columns of $A$ by rows of $B$:
    \begin{equation}
AB = \sum_{k=1}^m \begin{bmatrix} | \\ c_k^A \\ | \end{bmatrix} \begin{bmatrix} - & (l_k^B)^\trp & - \end{bmatrix} .
    \end{equation}
\end{enumerate}

\subsection{Special Matrices}
\begin{itemize}
    \item \textbf{Symmetric}: $A \in \R^{n \times n}$ is symmetric if $A^\trp = A$. Symmetric matrices possess real eigenvalues and a complete set of orthogonal eigenvectors. For any matrix $M \in \R^{n \times n}$, the matrix $M + M^\trp$ and the product $M M^\trp$ are always symmetric.
    \item \textbf{Anti-Symmetric}: $A$ is anti-symmetric if $A^\trp = -A$. They have purely imaginary eigenvalues, and their diagonal elements are zero.
    \item \textbf{Orthogonal}: $A$ is orthogonal if $A^\trp A = A A^\trp = I$. Orthogonal transformations preserve vector lengths (norms) and angles. An iconic example is the rotation matrix in $\R^2$:
    \begin{equation}
R(\theta) = \begin{bmatrix} \cos\theta & -\sin\theta \\ \sin\theta & \cos\theta \end{bmatrix} .
    \end{equation}
\end{itemize}

\subsection{Trace and Determinant Operators}
\begin{itemize}
    \item \textbf{Trace}: The trace of a square matrix $A \in \R^{n \times n}$ is the sum of its diagonal entries, which also equals the sum of its eigenvalues:
    \begin{equation}
\tr[A] = \sum_{i=1}^n a_{ii} = \sum_{i=1}^n \lambda_i .
    \end{equation}
    The trace is linear, $\tr[\alpha A + \beta B] = \alpha\tr[A] + \beta\tr[B]$, and cyclic, $\tr[AB] = \tr[BA]$.
    \item \textbf{Determinant}: The determinant is the product of its eigenvalues:
    \begin{equation}
\det[A] = \prod_{i=1}^n \lambda_i .
    \end{equation}
    It satisfies the multiplicative property: $\det[AB] = \det[A]\det[B]$.
\end{itemize}

\subsection{Eigenvalues and Spectral Decomposition}
For a square matrix $A \in \R^{n \times n}$, a non-zero vector $v \in \mathbb{C}^n$ is an eigenvector corresponding to eigenvalue $\lambda \in \mathbb{C}$ if:
\begin{equation}
A v = \lambda v .
\end{equation}
If $A$ is diagonalizable, it possesses $n$ linearly independent eigenvectors. Letting $V = [v_1 \dots v_n]$ be the modal matrix and $\Lambda = \diag(\lambda_1 \dots \lambda_n)$ the diagonal eigenvalue matrix, we have the \textbf{spectral decomposition}:
\begin{equation}
A = V \Lambda V^{-1} .
\end{equation}
This decomposition provides powerful computational simplifications:
\begin{itemize}
    \item \textbf{Matrix Inverse}: $A^{-1} = V \Lambda^{-1} V^{-1} = V \diag(\frac{1}{\lambda_1} \dots \frac{1}{\lambda_n}) V^{-1}$.
    \item \textbf{Matrix Powers}: $A^k = V \Lambda^k V^{-1} = V \diag(\lambda_1^k \dots \lambda_n^k) V^{-1}$.
    \item \textbf{Matrix Exponential}: $e^A = V e^\Lambda V^{-1} = V \diag(e^{\lambda_1} \dots e^{\lambda_n}) V^{-1}$.
\end{itemize}

\section{Solution of Continuous-Time LTI State-Space Models}
Consider the unforced (autonomous) LTI state-space equation:
\begin{equation}
\dot{x}(t) = A x(t), \quad x(0) = x_0 .
\end{equation}
By analogy with scalar differential equations, the solution is written in terms of the \textbf{matrix exponential}:
\begin{equation}
x(t) = e^{At} x_0 .
\end{equation}
\begin{mydefinition}{Matrix Exponential}
For any square matrix $A \in \R^{n \times n}$ and scalar $t$, the matrix exponential is defined by the globally convergent power series:
\begin{equation}
e^{At} = I + At + \frac{1}{2!} A^2 t^2 + \frac{1}{3!} A^3 t^3 + \dots = \sum_{k=0}^\infty \frac{1}{k!} A^k t^k .
\end{equation}
\end{mydefinition}

Calculating the matrix exponential numerically is a highly non-trivial task. In their landmark paper \textit{"Nineteen Dubious Ways to Compute the Exponential of a Matrix"} (SIAM Review, 1978 and 2003 updates), Cleve Moler and Charles Van Loan explored the mathematical and numerical properties of 19 different methods (including series approximation, ODE solvers, polynomial methods, and matrix decomposition), highlighting issues of numerical stability, stiffness, and floating-point errors.

\section{From Ordinary Differential Equations to State-Space Representations}
An $n$-th order linear ordinary differential equation can be transformed into an equivalent first-order vector-matrix state-space model.

Consider the scalar differential equation without input derivatives:
\begin{equation}
y^{(n)}(t) + a_{n-1} y^{(n-1)}(t) + \dots + a_1 \dot{y}(t) + a_0 y(t) = u(t) .
\end{equation}
To define a state-space model, we define the state variables as the output and its first $n-1$ derivatives:
\begin{equation}
x_1 = y(t), \quad x_2 = \dot{y}(t), \quad x_3 = \ddot{y}(t), \quad \dots \quad x_n = y^{(n-1)}(t) .
\end{equation}
Differentiating these state variables yields the relations:
\begin{align}
    \dot{x}_1 &= x_2 \\
    \dot{x}_2 &= x_3 \\
    \vdots \\
    \dot{x}_{n-1} &= x_n \\
\dot{x}_n &= y^{(n)}(t) = -a_0 x_1 - a_1 x_2 - \dots - a_{n-1} x_n + u(t) .
\end{align}
Expressing these relations in matrix-vector form gives the \textbf{Controllability Canonical Form}:
\begin{equation}
    \begin{bmatrix} \dot{x}_1 \\ \dot{x}_2 \\ \vdots \\ \dot{x}_{n-1} \\ \dot{x}_n \end{bmatrix} = \begin{bmatrix}
        0 & 1 & 0 & \cdots & 0 \\
        0 & 0 & 1 & \cdots & 0 \\
        \vdots & \vdots & \vdots & \ddots & \vdots \\
        0 & 0 & 0 & \cdots & 1 \\
        -a_0 & -a_1 & -a_2 & \cdots & -a_{n-1}
\end{bmatrix} \begin{bmatrix} x_1 \\ x_2 \\ \vdots \\ x_{n-1} \\ x_n \end{bmatrix} + \begin{bmatrix} 0 \\ 0 \\ \vdots \\ 0 \\ 1 \end{bmatrix} u(t) .
\end{equation}
The output measurement equation is:
\begin{equation}
y(t) = \begin{bmatrix} 1 & 0 & 0 & \cdots & 0 \end{bmatrix} \begin{bmatrix} x_1 \\ x_2 \\ \vdots \\ x_{n-1} \\ x_n \end{bmatrix} .
\end{equation}
This standard representation is extremely powerful for feedback controller design, as the coefficients of the system's characteristic polynomial appear directly in the last row of the dynamic matrix.
